\section*{Hoofdstuk \ref{ch:b2:precambrian-earth} --- De precambrische Aarde}

\begin{solution}{exo:b2:precambrian-earth:1}
(a) \omterm{def:b2:precambrian-earth:eons}{Hadeïcum} $4567 - 4031 = \qty{536}{Myr}$; \omterm{def:b2:precambrian-earth:eons}{Archaïcum} $4031 - 2500 = \qty{1531}{Myr}$;
\omterm{def:b2:precambrian-earth:eons}{Proterozoïcum} $2500 - 538.8 = \qty{1961}{Myr}$. (b) $(4567 - 538.8)/4567 = 0.882$: ongeveer zeven
achtste. (c) $4567 - 4404 = \qty{163}{Myr}$ voor de zirkoon; $4567 - 4030 = \qty{537}{Myr}$ voor
de oudste gesteenten.
\end{solution}

\begin{solution}{exo:b2:precambrian-earth:2}
(a) De factoren $e^{\lambda t}$ met $\lambda = \ln2/T$ en $t = \qty{3}{Gyr}$ zijn 1.593
(\ce{^{238}U}), 19.18 (\ce{^{235}U}), 1.160 (\ce{^{232}Th}) en 5.292 (\ce{^{40}K}): de
bijdragen waren 3.053, 1.573, 2.424 en \qty{4.329}{pW}, \qty{11.38}{pW} in totaal, tegen
\qty{4.906}{pW} vandaag. (b) $11.38/4.906 = 2.32$. (c) Thorium-232 leidt vandaag; kalium-40 leidde
toen.
\end{solution}

\begin{solution}{exo:b2:precambrian-earth:3}
(a) $1616 + 18.0\times(28 - 32.2) = \qty{1540}{\celsius}$. (b) $32.2 + (1500 - 1616)/18.0 =
\qty{25.8}{\percent}$. (c) Een analyse van het totale gesteente is niet de samenstelling van een vloeistof: een
komatiietstroom kan olivijnkristallen bevatten die uit de vloeistof zijn opgehoopt en MgO toevoegen. Zijn
\qty{38}{\percent} ligt voorbij de \qty{32}{\percent} van de rijkste komatiietvloeistoffen, en
voorbij de experimenten waarop de lijn berust.
\end{solution}

\begin{solution}{exo:b2:precambrian-earth:4}
Oudste gesteenten, \qty{4.03}{Ga}; eerste stromatolieten, \qty{3.43}{Ga}; einde van de koepels van de Pilbara,
\qty{3.1}{Ga}; \omterm{def:b2:precambrian-earth:goe}{Grote Oxidatie}, ongeveer 2.4 tot \qty{2.1}{Ga}; laatste gewone \omterm{def:b2:precambrian-earth:greenstone}{komatiieten},
\qty{1.9}{Ga}; tweede zuurstofstijging, ongeveer \qty{0.8}{Ga}; einde van de Marinoische ijstijd,
\qty{0.635}{Ga}; eerste fossielen van het Ediacarium, ongeveer \qty{0.57}{Ga}.
\end{solution}

\begin{solution}{exo:b2:precambrian-earth:5}
(a) $\tfrac14\times1361\times0.706 = \qty{240.2}{W.m^{-2}}$ geabsorbeerd en $\sigma\times288^4 =
\qty{390.1}{W.m^{-2}}$, dus $\varepsilon = 240.2/390.1 = 0.616$. (b) $\tfrac14\times1361\times0.4
= \qty{136.1}{W.m^{-2}}$, $T^4 = 136.1/(0.616\times\num{5.670e-8})$, $T = \qty{249.8}{K}$.
(c) Bij \qty{275.5}{K} is het albedo $0.6 - 0.306\times12.5/20 = 0.4088$: $R = 340.25\times0.5912 =
\qty{201.2}{W.m^{-2}}$ en $E = 0.616\,\sigma\times275.5^4 = \qty{201.2}{W.m^{-2}}$. (d) $R' =
340.25\times0.306/20 = \qty{5.21}{W.m^{-2}.K^{-1}}$, $E' = 4\varepsilon\sigma T^3 =
\qty{2.92}{W.m^{-2}.K^{-1}}$: $R' > E'$, dus het evenwicht is instabiel
(\cref{prop:b2:precambrian-earth:ebm}).
\end{solution}

\begin{solution}{exo:b2:precambrian-earth:6}
(a) $1/(1 + 0.4\times2.4/4.567) = 0.826$; $1/(1 + 0.4\times0.65/4.567) = 0.946$. (b) De warme
toestand vraagt minstens 0.932 van het huidige zonlicht: onmogelijk \qty{2.4}{Ga} geleden, mogelijk
\qty{0.65}{Ga} geleden. (c) $0.826/0.946 = 0.873$: \qty{13}{\percent} zwakker, dicht bij de
geciteerde \qty{10}{\percent}, die uit een ander zonnemodel komt.
\end{solution}

\begin{solution}{exo:b2:precambrian-earth:7}
(a) Sturtische ijstijd $716.9 - 658.9 = \qty{58.0}{Myr}$; Marinoische $644.5 - 635.4 = \qty{9.1}{Myr}$;
interglaciaal $658.9 - 644.5 = \qty{14.4}{Myr}$. (b) $(58.0 + 9.1)/85 = 0.79$: ongeveer vier vijfde
van de periode. (c) Haar begin ligt alleen ingesloten tussen 649.9 en \qty{639.0}{Ma}: met deze
grenzen duurde de Marinoische ijstijd tussen ongeveer 3 en \qty{15}{Myr}.
\end{solution}

\begin{solution}{exo:b2:precambrian-earth:8}
(a) $\tau = \frac{0.05}{0.95}\times\frac{44.01}{28.97}\times\frac{\num{5.1e18}}{\num{0.26e12}} =
0.05263\times1.519\times\num{1.962e7} = \qty{1.57}{Myr}$. (b) $0.05/\num{278e-6} = 180$ keer.
(c) Twee keer zo lang, \qty{3.14}{Myr}.
\end{solution}

\begin{solution}{exo:b2:precambrian-earth:9}
(a) De stenen vielen in zachte modder en bogen de laminae ervan omlaag: het zijn valstenen, losgelaten door
drijvend ijs. (b) IJs dat de zee bereikt op \ang{10} van de evenaar is het verwachte spoor van een
sneeuwbal, en zeer moeilijk te verkrijgen als alleen de poolgebieden bevroren zijn. (c) Een \omterm{def:b2:precambrian-earth:snowball}{kapcarbonaat}:
het legt het plotselinge einde van de ijstijd vast, met intense verwering en de overstroming van de
plateaus door het smeltwater.
\end{solution}

\begin{solution}{exo:b2:precambrian-earth:10}
(a) $0.94 > 0.932$, het warme kantelpunt: het warme evenwicht, $288\times0.94^{1/4} =
\qty{283.6}{K}$, ligt net boven \qty{283}{K}; en $0.94 < 1.227$, dus ook de sneeuwbal bestaat.
(b) $249.8\times0.94^{1/4} = \qty{246.0}{K}$. (c) $\varepsilon = 0.94\times1361\times0.4/(4\sigma
\times263^4) = 511.7/1085.1 = 0.472$. (d) $\sigma\times263^4 = \qty{271.3}{W.m^{-2}}$, en $(0.616
- 0.472)\times271.3 = \qty{39}{W.m^{-2}}$.
\end{solution}

\begin{solution}{exo:b2:precambrian-earth:11}
(a) $S_{\text{m}}/S_{\text{f}} = (T_1/T_2)^4\,(1 - a_{\text{w}})/(1 - a_{\text{i}})$: de
emissiviteit valt weg. $(263/283)^4\times0.706/0.4 = 0.7459\times1.765 = 1.317$. (b) Een kleiner
albedocontrast tussen ijs en open water, zoals donkerder, stoffig ijs, of een bredere
temperatuurhelling. (c) Een sneeuwbal weerkaatst zoveel zonlicht dat hij koud blijft bij elk zonlicht of
broeikaseffect waarbij een warme Aarde zou bevriezen; om hem te ontdooien moet de forcering die van het
bevriezen overtreffen met de volle breedte van de lus.
\end{solution}

\begin{solution}{exo:b2:precambrian-earth:12}
In de buurt van $T^*$ is $C\dot\delta = -(E' - R')\delta$, dus $\delta \propto e^{-t/\tau}$ met $\tau =
C/(E' - R')$. In het warme evenwicht is $R' = 0$ en $E' = 4\times0.616\times\num{5.670e-8}\times
288^3 = \qty{3.34}{W.m^{-2}.K^{-1}}$; $C = 1025\times3992\times50 = \qty{2.05e8}{J.m^{-2}.K^{-1}}$,
dus $\tau = \qty{6.1e7}{s}$, ongeveer \qty{1.9}{yr}. In dit model houdt het warme evenwicht $R' =
0$ tot het \qty{283}{K} bereikt, dus $\tau$ blijft tot aan het kantelpunt dicht bij twee jaar, waarna
de toestand verdwijnt. Met een geleidelijk variërend albedo groeit $R'$ naarmate het evenwicht de
plooi nadert, waar $E' = R'$: $\tau$ groeit onbegrensd, een vertraging die het kantelpunt
aankondigt.
\end{solution}

\begin{solution}{pb:b2:precambrian-earth:1}
\textbf{1.} $\num{0.0203e-6}\times0.9929\times\num{95.13e-6} = \qty{1.92}{pW.kg^{-1}}$
(\ce{^{238}U}); $\num{0.0203e-6}\times0.00711\times\num{568.47e-6} = \qty{0.0820}{pW.kg^{-1}}$
(\ce{^{235}U}); $\num{0.0795e-6}\times\num{26.28e-6} = \qty{2.09}{pW.kg^{-1}}$ (\ce{^{232}Th});
$\num{240e-6}\times\num{1.20e-4}\times\num{28.47e-6} = \qty{0.820}{pW.kg^{-1}}$ (\ce{^{40}K}).
\textbf{2.} \qty{4.91}{pW.kg^{-1}}.
\textbf{3.} $\ln2/T$: 0.155, 0.985, 0.0495 en \qty{0.555}{Gyr^{-1}}.
\textbf{4.} Factoren 1.52, 14.3, 1.14 en 4.48: 2.91, 1.17, 2.39 en \qty{3.67}{pW.kg^{-1}}.
\textbf{5.} \qty{10.1}{pW.kg^{-1}}, 2.07 keer het huidige.
\textbf{6.} Kalium-40. Zijn halveringstijd, \qty{1.25}{Gyr}, is kort: sindsdien zijn er meer dan twee halveringstijden
verstreken, en meer dan drie kwart van het \ce{^{40}K} van toen is vervallen.
\textbf{7.} $1/(1 + 0.4\times2.7/4.567) = 0.809$; $0.809\times1361 = \qty{1101}{W.m^{-2}}$.
\textbf{8.} $288\times0.809^{1/4} = \qty{273}{K}$, onder \qty{283}{K}: geen geldige warme toestand;
de Aarde zou bevriezen.
\textbf{9.} $(1101\times0.4/(4\times0.616\,\sigma))^{1/4} = \qty{237}{K}$.
\textbf{10.} $0.616\times0.809 = 0.498$.
\textbf{11.} $0.004/\num{278e-6} = 14$ tot $0.03/\num{278e-6} = 108$ keer.
\textbf{12.} De thermostaat van de silicaatverwering: de verwering vertraagt in een koud klimaat, zodat
vulkanische koolstofdioxide zich ophoopt, en versnelt in een warm klimaat.
\textbf{13.} 0.826 van het huidige, onder 0.932: geen warme toestand; de Aarde met de huidige lucht zou een
sneeuwbal zijn.
\textbf{14.} $\frac{0.1}{0.9}\times1.519\times\num{1.962e7} = \qty{3.31}{Myr}$.
\textbf{15.} $\frac{0.2}{0.8}\times1.519\times\num{1.962e7} = \qty{7.45}{Myr}$.
\textbf{16.} Het \omterm{def:b2:precambrian-earth:snowball}{kapcarbonaat}, een gelamineerde dolosteen die scherp op de glaciale afzettingen ligt,
neergeslagen uit een zee die gevoed werd door intense verwering in de hete nasleep en overstroomd door
smeltwater.
\textbf{17.} IJzer(II) blijft alleen in water zonder zuurstof opgelost: het ijzer hoopt zich op in de
anoxische oceaan onder het ijs en slaat neer als ijzer(III) waar zuurstof aankomt, bijvoorbeeld met
smeltwater.
\textbf{18.} $(1616 - 1450)/(32.2 - 23) = \qty{18.0}{\celsius}$ per procent.
\textbf{19.} $1616 + 18.0\times(30 - 32.2) = \qty{1576}{\celsius}$.
\textbf{20.} \qty{1616}{\celsius}, het gemeten punt.
\textbf{21.} $1616 - 18.0\times20.2 = \qty{1252}{\celsius}$.
\textbf{22.} $1616 - 20\times20.2 = \qty{1212}{\celsius}$: \qty{40}{\celsius} lager, een tiende van
het verschil tussen de lava's; de conclusie hangt niet van de keuze af.
\textbf{23.} \omterm{def:b2:precambrian-earth:greenstone}{Komatiieten} kwamen uit pluimen, heter dan de omringende mantel, en uit hoge smeltgraden
diep beneden; de temperatuur van een lava is niet de potentiële temperatuur van haar bron. De
mantel als geheel was 200 tot \qty{300}{K} warmer.
\textbf{24.} De heetste komatiietlava was ongeveer \textbf{\qty{364}{\celsius}} heter dan een
primair basalt van de huidige ruggen.
\end{solution}
